\documentclass{article}
\usepackage{graphicx}
\usepackage[margin=1in]{geometry}
\usepackage{enumitem}
\usepackage{booktabs}
\usepackage{array}
\usepackage{parskip}
\usepackage{hyperref}
\setlist{noitemsep, topsep=3pt, leftmargin=*}

\title{Beyond the Shadow of a Doubt: Benchmarking MCMC Sampling Algorithms
  for Multimodal Astrophysical Inference\\[0.5em]
  \large Progress Report}
\author{Lucas Arthur \and Samson Mercier}
\date{April 13, 2026}

\begin{document}

\maketitle

\begin{abstract}
As astrophysics enters the era of petabyte-scale datasets, the choice of
inference algorithm and implementation for a research problem could make the
difference between an intractable or incorrect analysis and a major scientific
discovery. The high dimensionality and multi-modality of posterior distributions
can pose problems for popular inference tools. To ensure that the astrophysics
community is fully leveraging available inference algorithms, we benchmark a
set of widespread tools from across the astrophysical literature. We benchmark
seven inference algorithms---random walk Metropolis-Hastings (RWMH),
affine-invariant ensemble sampling, parallel tempering, NUTS, SMC with
tempering, non-reversible parallel tempering, and nested sampling (as a
non-MCMC comparator)---on two test problems: a synthetic
8-component 2D Gaussian mixture with known ground truth, and a simulated
exoplanet transit with stellar spots. All experiments run on AMD EPYC 9474F
48-core CPUs and Nvidia L40S GPUs using JAX to ensure
implementation-agnostic comparison. We evaluate each sampler on mode
recovery, posterior accuracy, and computational efficiency (ESS per log-likelihood evaluation and ESS per
wall-clock second), and will produce a concrete set of recommendations for practitioners.
\end{abstract}

% -----------------------------------------------------------------------
\section{Project Overview}

The central question of this project is: \emph{Which parallelizable MCMC-family samplers most reliably recover multi-modal posteriors in astrophysical settings, and at what computational cost?} We aim to benchmark seven algorithms: Random-Walk Metropolis Hastings (RWMH), affine-invariant ensemble sampling \cite{goodman2010,emcee}, reversible parallel tempering (Stochastic Even Odd, SEO) \cite{vousden2016}, non-reversible parallel tempering (Deterministic Even Odd, DEO) \cite{syed2019}, NUTS, SMC with tempering \cite{delmoral2006}, and nested sampling \cite{speagle2020} as a non-MCMC reference.\footnote{Ensemble Kalman sampling \cite{garbuno2020} was originally planned as an eighth algorithm but is deferred pending completion of all other sampler runs; it will be included if time permits.}

We evaluate each sampler on two test problems. The first is an 8-component 2D Gaussian mixture with fully known ground-truth modes, weights, and covariances, enabling precise measurement of mode recovery and posterior accuracy. This problem is advantageous for two reasons: its simple analytical form and low-dimension provides an intuitive visualization of each sampler's results, yet, despite this simplicity, accurately recovering all eight components remains a challenging task. The second is a simulated exoplanet transit with active regions on the host star's surface, a problem known to produce multi-modal and highly correlated posteriors \cite{Pinhas2018, Rackham2018, Sagynbayeva2025}. Active regions have been a growing source of bias in exoplanet science \cite{Murray2025}, yet no systematic benchmarking has identified the most reliable sampler for jointly constraining active region and exoplanet properties. The evaluation criteria are discussed in Section~6.

The full implementation stack is \texttt{JAX}\footnote{\url{https://github.com/jax-ml/jax}}~\cite{jax2018github} / \texttt{BlackJAX}\footnote{\url{https://github.com/blackjax-devs/blackjax}}~\cite{blackjax} /
\texttt{NumPyro}\footnote{\url{https://github.com/pyro-ppl/numpyro}}~\cite{numpyro}, ensuring hardware-accelerated, implementation-agnostic
comparison across all algorithms.

% -----------------------------------------------------------------------
\section{Division of Labor Update}

The division of labor has been reshuffled relative to the original proposal.
Lucas Arthur has taken ownership of both parallel tempering variants---reversible
SEO and non-reversible DEO---in addition to his originally assigned algorithms
(NUTS and SMC with tempering).

In addition to his originally assigned algorithms (RWMH,
affine-invariant MCMC, and nested sampling), Samson Mercier has taken on minor development of the forward transit model.

The ensemble Kalman sampler (EKS) is de-prioritized. Implementation is on hold
until all other samplers have been run and results collated across both test
problems. Samson Mercier will implement it if time permits; it will be omitted
from the final comparison if it is not ready in time.

% -----------------------------------------------------------------------
\section{Progress to Date}

\subsection{Gaussian Mixture Model --- Substantially Complete}

Seven of the eight samplers have been fully implemented and run on the Gaussian
mixture problem: RWMH, affine-invariant MCMC, reversible PT (SEO),
non-reversible PT (DEO), NUTS, SMC with tempering, and nested sampling. EKS is
the sole missing algorithm; it is deferred pending completion of transit model
runs.

Each sampler produces the following outputs: posterior samples in a
standardized format, corner plots, and a diagnostics record containing ESS,
$\hat{R}$, acceptance rates, recovered mode weights, and inter-mode transition
counts (for chain-based methods).

\subsection{Stellar Spot Transit Model --- Forward Model Complete, Sampling Not Yet Started}

The transit forward model has been implemented using the \texttt{SAJAX} library\footnote{\url{https://github.com/SamMerc/sajax}}
with a NumPyro likelihood and BlackJAX-compatible log-density interface.

Each active region is described by four parameters (longitude, latitude, angular size, and contrast); three stellar parameters (two limb-darkening coefficients and rotational period) and seven planetary orbital parameters complete the model, yielding 17 free parameters after fixing the mid-transit time (typically known to minute-level precision). Tight Gaussian priors are placed on the orbital and rotation periods, centered on truth.

Samson Mercier has validated the forward model against independent transit-fitting codes. Sampler scripts for the transit problem have not yet been written. Writing and running these scripts is the primary remaining task.

\begin{figure}[h]
\centering
\fbox{\begin{minipage}[c][5cm][c]{0.75\textwidth}
\centering\textit{[Placeholder: synthetic transit light curve]}
\end{minipage}}
\caption{Synthetic exoplanet transit light curve generated by the \texttt{SAJAX} forward
  model. The $x$-axis shows time from mid-transit (hours) and the $y$-axis shows
  normalized stellar flux. The transit profile is modulated by two active regions---a
  dark spot and a bright facula---on the stellar surface.}
\label{fig:lightcurve}
\end{figure}

% -----------------------------------------------------------------------
\section{Preliminary Results: Gaussian Mixture}

Notable observations include:

\begin{itemize}
  \item \textbf{RWMH} shows poor mixing between modes. The chain rarely crosses
    the low-density region separating mixture components, resulting in
    substantially fewer inter-mode transitions and lower ESS per evaluation than
    all other methods. Performance is sensitive to initialization; chains started
    near a single mode frequently fail to discover remaining components.
  \item \textbf{Affine-invariant ensemble sampling} shows moderate inter-mode
    mixing when walkers are initialized broadly across the support, but degrades
    substantially when the ensemble is concentrated near a single mode.
  \item \textbf{Reversible PT (SEO)} achieves inter-mode mixing through
    temperature swaps, though at a lower rate than DEO at matched compute budgets.
  \item \textbf{Non-reversible PT (DEO)} achieves substantially more inter-mode
    transitions than reversible PT (SEO) at matched compute budgets, consistent
    with the theoretical prediction that non-reversible temperature swaps improve
    round-trip rates \cite{syed2019}.
  \item \textbf{Nested sampling} reliably recovers all eight modes. Its ESS per
    log-density evaluation is lower than gradient-based methods, but it provides
    an accurate evidence estimate and serves as a robust reference.
  \item \textbf{NUTS} and \textbf{SMC with tempering} show strong ESS per
    evaluation on this low-dimensional problem; their relative performance on the
    higher-dimensional transit model remains to be seen.
\end{itemize}
Initialization sensitivity is a recurring theme across chain-based methods; we
are actively exploring parameter tuning to establish fair comparison baselines
ahead of the transit model runs.

\begin{figure}[h]
\centering
\fbox{\begin{minipage}[c][7cm][c]{0.75\textwidth}
\centering\textit{[Placeholder: SMC corner plot]}
\end{minipage}}
\caption{Corner plot of the SMC with tempering posterior on the 8-component 2D Gaussian
  mixture. Diagonal panels show 1D marginal posteriors with parameter value on the
  $x$-axis and normalized density on the $y$-axis; the off-diagonal panel shows the 2D
  joint marginal with $x_1$ on the $x$-axis and $x_2$ on the $y$-axis. Black crosses
  mark the true component means.}
\label{fig:smc-corner}
\end{figure}

% -----------------------------------------------------------------------
\section{Revised Plan and Remaining Tasks}

Table~\ref{tab:plan} lists the remaining tasks and internal deadlines.

\begin{table}[h]
\centering\small
\begin{tabular}{@{}p{0.60\textwidth}p{0.14\textwidth}p{0.10\textwidth}@{}}
\toprule
\textbf{Task} & \textbf{Owner} & \textbf{Target} \\
\midrule
Run all 7 implemented samplers on transit model & Both   & Apr 25 \\
Collate all results; build comparison tables    & Both   & Apr 25 \\
Implement EKS (if time permits)                 & Samson & Apr 25 \\
Produce final figures                           & Both   & Apr 28 \\
Write final report                              & Both   & May 1  \\
\bottomrule
\end{tabular}
\caption{Remaining tasks and internal deadlines.}
\label{tab:plan}
\end{table}

% -----------------------------------------------------------------------
\section{Planned Figures and Evaluation Framework}

Comparing the seven algorithms requires a unified evaluation framework because they produce
fundamentally different outputs: correlated samples (RWMH, PT, NUTS), weighted
particle sets (SMC, nested sampling), or continuously evolved ensembles (affine-invariant).
We therefore report efficiency as ESS per log-likelihood evaluation and ESS per
wall-clock second. The primary evaluation criteria are:
\begin{enumerate}
    \item \textbf{Mode-recovery metrics:} the per-mode weight error $|\hat{w}_k - w_k|$,
      total variation distance from the true mixture, and (for chain methods) inter-mode
      transition count. For the Gaussian mixture, these are evaluated directly against
      ground truth; mode recovery is excluded from the transit model comparison since
      ground-truth modes are not readily available.
    \item \textbf{Computational efficiency:} ESS per log-likelihood evaluation and ESS
      per wall-clock second. For chain-based methods, warm-up evaluations and warm-up
      time are included in the respective denominators.
\end{enumerate}
Secondary diagnostics include the Gelman--Rubin potential scale reduction factor
($\hat{R}$), acceptance rates, and---for parallel tempering variants---swap acceptance
rates and round-trip counts.

The main body will contain three figures:
\begin{enumerate}
    \item \textbf{KDE posterior overlay (Gaussian mixture):} one panel per sampler, arranged
      in a grid. The $x$-axis is the first mixture dimension $x_1$ and the $y$-axis is the
      second dimension $x_2$; filled contours show the true mixture density and overlaid KDE
      contours show the sampler's recovered posterior.
    \item \textbf{Transit model comparison:} a two-panel corner plot contrasting the best-
      and worst-performing samplers on the transit problem. Along the diagonal, each panel
      shows a 1D marginal posterior with the parameter value on the $x$-axis and normalized
      density on the $y$-axis. Off-diagonal panels show 2D joint marginals with the two
      respective parameter values on the $x$- and $y$-axes, with contours at the 68th and
      95th percentile levels.
    \item \textbf{Summary table:} rows are samplers; columns are ESS per log-likelihood
      evaluation, ESS per wall-clock second (both reported for each test problem),
      fraction of modes recovered (Gaussian mixture only), $\hat{R}$, and acceptance rate.
\end{enumerate}

Full per-sampler diagnostics are relegated to an appendix. For each sampler on both test
problems, the appendix includes: a scatter plot of raw posterior samples (Gaussian mixture:
$x$-axis $x_1$, $y$-axis $x_2$; transit model: one 2D projection per parameter pair),
trace plots for each parameter ($x$-axis: iteration number; $y$-axis: parameter value),
and a full corner plot (diagonal panels: parameter value vs.\ normalized density; off-diagonal
panels: pairwise parameter values with density contours).

\bibliographystyle{unsrt}
\bibliography{refs}

\end{document}
